\chapter{Точная LCF-проверка двухсекторной неподвижной алгебры}
\label{ch:v8-markov-fixed-algebra-lcf-migration-gate}

\section{От численного ядра к точному коммутанту}

Исходный гейт получил двумерную неподвижную алгебру диагонализацией
генератора на $221$-мерном пространстве наблюдаемых. Перепроверим результат
без численного допуска.

Коммутант полного калибровочного представления внутри двух углов концевых секторов имеет
размерность
\begin{equation}
 \bigl(1^2+2^2+1^2\bigr)_{E_s}
 +\bigl(1^2+1^2+2^2+1^2\bigr)_{E_t}=13.
 \label{eq:v8-lcf-gauge-commutant-thirteen}
\end{equation}
Это следует из точных изотипических кратностей исходных  и целевых модулей.

Пусть $X_s\oplus X_t$ --- общий элемент этого 13-мерного коммутанта, а
$A:\mathbb C^{11}\to\mathbb C^{10}$ --- физическая матрица инцидентности с
элементами $0$ и $1$. Коммутирование с полным самосопряжённым оператором
\begin{equation}
 D_A=\begin{pmatrix}0&A^*\\A&0\end{pmatrix}
 \end{equation}
эквивалентно двум условиям
\begin{equation}
 AX_s=X_tA,
 \qquad
 X_sA^*=A^*X_t.
 \label{eq:v8-lcf-two-sided-linking-commutant}
\end{equation}

\section{Почему необходимы обе половины}

Одностороннее условие $AX_s=X_tA$ имеет четырёхмерное ядро. Поэтому оно не
доказывает двухсекторный результат. Совместная точная система имеет форму
$220\times13$, ранг $11$ и нульмерность
\begin{equation}
 13-11=2.
 \label{eq:v8-lcf-fixed-nullity-two}
\end{equation}
Её базис состоит ровно из цветного и бесцветного проекторов:
\begin{equation}
 \ker=\mathbb CP_q\oplus\mathbb CP_\ell,
 \qquad
 \rank P_q=12,
 \qquad
 \rank P_\ell=9.
 \label{eq:v8-lcf-fixed-projector-ranks}
\end{equation}
LCF-ядро отдельно проверяет идемпотентность, ортогональность и равенство
$P_q+P_\ell=I_{21}$.

\begin{theorem}[Точная двухсекторная неподвижная алгебра]
Общий коммутант полного калибровочного представления концевых секторов и физического
самосопряжённого связывающего оператора точно равен
$\mathbb CP_q\oplus\mathbb CP_\ell$. Его размерность равна двум, а ранги
дополнительных проекторов равны $12$ и $9$. Численный допуск не используется.
\end{theorem}

\begin{nogocriterion}[Запрет одностороннего сертификата]
Нельзя доказывать неподвижную алгебру только уравнением $AX_s=X_tA$: в
данной модели оно оставляет четыре направления. Требуется полный коммутант
$D_A$, включающий сопряжённое уравнение.
\end{nogocriterion}

\begin{protocol}
\begin{enumerate}
 \item калибровочный коммутант углов концевых секторов: \textbf{$13$};
 \item одностороннее связывающее ядро: \textbf{$4$};
 \item полная система: \textbf{$220\times13$};
 \item точный ранг системы: \textbf{$11$};
 \item полное неподвижное ядро: \textbf{$2$};
 \item ранги проекторов: \textbf{$12+9$};
 \item совпадение с прежним численным гейтом: \textbf{полное};
 \item статус: \textbf{проверено LCF};
 \item следующий перенос: \textbf{полный построитель GKSL/QMS}.
\end{enumerate}
\end{protocol}

Машинный сертификат сохранён в
\path{s2t_v8_markov_fixed_algebra_lcf_migration_gate_results.json}.